\chapter{Összefoglalás és javaslatok}

Jelen fejezet összefoglalja a dolgozat főbb eredményeit, levon következtetéseket a
posztkvantum-kriptográfiai algoritmusok IoT alkalmazhatóságáról, és javaslatokat
fogalmaz meg mind a gyakorlati bevezetés, mind a jövőbeli kutatások irányára vonatkozóan.

%----------------------------------------------------------------------------
\section{A dolgozat eredményeinek összefoglalása}
%----------------------------------------------------------------------------

% TODO: A CRYSTALS-Kyber és BIKE algoritmusok működésének, matematikai alapjainak
% és biztonsági feltevéseinek összefoglalása; mi az, amit a dolgozat hozzátesz
% a meglévő irodalomhoz.

% TODO: A legfontosabb mérési számok egy bekezdésben (pl. „A Kyber-512 encapsulation
% X ms alatt futott le, ami Y-szor gyorsabb, mint a BIKE-L1 megfelelője...").

%----------------------------------------------------------------------------
\section{Javaslatok IoT rendszerek PQC migrációjához}
%----------------------------------------------------------------------------

% TODO: Mikor válasszuk a Kybert, mikor a BIKE-ot — döntési szempontok (sebesség,
% kulcsméret, energiaigény, szabványosítási státusz).

% TODO: Klasszikus + PQC KEM kombináció ajánlása az átmeneti időszakra;
% hogyan valósítható meg wolfSSL-lel / OpenSSL-OQS-sel.

% TODO: Hogyan tervezzük meg az IoT firmware-t, hogy az algoritmusok cserélhetők
% legyenek (pl. absztrakciós réteg, szekvenciális frissítési stratégia).

% TODO: DTLS 1.3 + CoAP használata MQTT helyett erősen korlátozott eszközökön;
% tanúsítványlánc-méret optimalizálása; brókeroldali gyorsítás lehetőségei.

%----------------------------------------------------------------------------
\section{Jövőbeli kutatási irányok}
%----------------------------------------------------------------------------

% TODO: Mérések valódi Cortex-M0/M3 eszközön; ARM TrustZone / HSM integrációja.

% TODO: Valódi LoRaWAN vagy NB-IoT hálózati réteg bevonása; e2e késleltetés és
% csomagveszteség hatása PQC kézfogásokra.

% TODO: CRYSTALS-Dilithium alapú tanúsítványlánc vizsgálata IoT eszközökön;
% OCSP / CRL alternatívák kis sávszélességű csatornákhoz.

% TODO: Hardveres NTT-gyorsítás (pl. FPGA implementáció), alvó mód és
% ritka kulcscsere kombinálása az energiafogyasztás csökkentésére.
